\documentclass[12pt,a4paper]{article}

\usepackage{fontspec}
\setmainfont{Latin Modern Roman}
\usepackage[french]{babel}
\usepackage[a4paper,margin=2cm]{geometry}
\usepackage{setspace}
\usepackage{booktabs}
\usepackage{tabularx}
\usepackage{array}
\usepackage{enumitem}
\usepackage{xcolor}
\usepackage{amsmath}
\usepackage{fancyhdr}
\usepackage{hyperref}
\usepackage[most]{tcolorbox}
\usepackage{titlesec}

\definecolor{mainblue}{HTML}{123A63}
\definecolor{accent}{HTML}{008C95}
\definecolor{lightblue}{HTML}{EDF4F8}
\definecolor{warning}{HTML}{9A5B00}

\hypersetup{colorlinks=true,linkcolor=mainblue,urlcolor=accent}
\setstretch{1.15}
\setlength{\parindent}{0pt}
\setlength{\parskip}{0.45em}
\setlist[itemize]{leftmargin=1.4em,itemsep=0.25em,topsep=0.25em}
\setlist[enumerate]{leftmargin=1.5em,itemsep=0.25em,topsep=0.25em}
\titleformat{\section}{\Large\bfseries\color{mainblue}}{}{0pt}{}[\color{accent}\titlerule]
\titleformat{\subsection}{\large\bfseries\color{accent}}{}{0pt}{}

\pagestyle{fancy}
\setlength{\headheight}{15pt}
\fancyhf{}
\lhead{Support technique de soutenance}
\rhead{Recommandation emploi--compétences}
\cfoot{\thepage}
\renewcommand{\headrulewidth}{0.4pt}

\newtcolorbox{oralbox}{
  colback=lightblue,
  colframe=mainblue,
  boxrule=0.7pt,
  arc=2pt,
  title=Formulation orale conseillée,
  fonttitle=\bfseries
}

\newtcolorbox{alertbox}{
  colback=orange!5,
  colframe=warning,
  boxrule=0.7pt,
  arc=2pt,
  title=Point de vigilance,
  fonttitle=\bfseries
}

\begin{document}

\begin{titlepage}
  \centering
  \vspace*{2.2cm}
  {\color{mainblue}\rule{\textwidth}{1.2pt}}\\[0.8cm]
  {\Huge\bfseries Données, fine-tuning et diagnostics\\[0.25cm]
  du système hybride\par}
  \vspace{0.5cm}
  {\Large Support technique pour la présentation PowerPoint\par}
  \vspace{0.8cm}
  {\color{accent}\rule{0.72\textwidth}{0.8pt}}\\[1.2cm]
  {\large Système hybride de recommandation emploi--compétences\par}
  \vfill
  \begin{tcolorbox}[colback=lightblue,colframe=mainblue,arc=2pt,width=0.86\textwidth]
  Ce document rassemble les éléments à expliquer au jury : transformation des données,
  niveaux NCF, preuve du fine-tuning, paramètres optimisés, distinction avec LoRA,
  calibration bayésienne, structure du graphe, dépendance aux hubs et limites de validité.
  \end{tcolorbox}
  \vfill
  {\large Juin 2026}
\end{titlepage}

\tableofcontents
\newpage

\section{Préparation et normalisation des données}

La normalisation transforme des offres et des profils hétérogènes en informations comparables.
Elle ne crée pas une information absente : elle nettoie, organise et code ce qui est disponible.

\subsection{Chaîne de traitement}

\begin{enumerate}
  \item \textbf{Nettoyer les textes :} suppression des espaces inutiles, valeurs vides explicites
  et résidus provenant du scraping.
  \item \textbf{Uniformiser les écritures :} par exemple, \og douala\fg{}, \og DOUALA\fg{}
  et \og Douala \fg{} sont ramenés à une modalité commune.
  \item \textbf{Séparer les valeurs multiples :} \og Python ; Excel / SQL\fg{} devient une
  liste de trois compétences.
  \item \textbf{Coder les qualifications :} les diplômes et niveaux d'études sont rapprochés
  des niveaux de la Nomenclature camerounaise des formations (NCF).
  \item \textbf{Construire le texte d'encodage :} le champ \texttt{text\_to\_embed} réunit
  le titre, les compétences, le secteur, le niveau, l'expérience et la description disponibles.
\end{enumerate}

\subsection{Exemple simple}

\begin{tabularx}{\textwidth}{>{\bfseries}p{0.22\textwidth}X}
\toprule
Donnée brute & \og Data analyst ; douala ; Bac+5 ; Python/Excel/SQL \fg{} \\
\midrule
Donnée normalisée & Métier : Data Analyst ; ville : Douala ; niveau NCF : 8 ;
compétences : Python, Excel et SQL. \\
\bottomrule
\end{tabularx}

Les données normalisées alimentent deux représentations complémentaires : la mémoire
vectorielle compare le sens global des textes ; le graphe crée des relations explicites telles
que \texttt{CANDIDAT--POSSEDE--COMPETENCE} et \texttt{OFFRE--REQUIERT--COMPETENCE}.

\begin{alertbox}
Le pipeline ne supprime pas automatiquement toutes les lignes répétées. Il audite les doublons
et conserve les observations ; il retire surtout les répétitions internes aux listes, par exemple
\texttt{[Python, Python, Excel]} devient \texttt{[Python, Excel]}.
\end{alertbox}

\begin{oralbox}
« La normalisation met les offres et les candidats dans un langage commun. Elle rend les champs
comparables avant l'encodage vectoriel et avant la création des relations du graphe. »
\end{oralbox}

\newpage
\section{Niveaux de la NCF utilisés par le système}

Les codes NCF ne sont pas des notes. Ils ordonnent les types de formation et facilitent la
comparaison entre le niveau d'un candidat et celui demandé par une offre.

\renewcommand{\arraystretch}{1.12}
\begin{tabularx}{\textwidth}{>{\centering\arraybackslash}p{0.08\textwidth}X>{\raggedright\arraybackslash}p{0.31\textwidth}}
\toprule
\textbf{Code} & \textbf{Niveau de formation} & \textbf{Exemples} \\
\midrule
1 & Formation sur le tas & Apprentissage pratique, attestation \\
2 & Formation courte en centre agréé & Formation généralement inférieure à un an \\
3 & Primaire et post-primaire & CEP, SAR/SM \\
4 & Secondaire premier cycle & BEPC, CAP, BEP \\
5 & Secondaire second cycle & Baccalauréat, agent technique \\
6 & Postsecondaire court, Bac+2 & BTS, DUT, HND, DEUG \\
7 & Supérieur premier cycle, Bac+3 & Licence, ingénieur des travaux \\
8 & Supérieur second cycle, Bac+4 ou Bac+5 & Master I et II, maîtrise, ingénieur de conception \\
9 & Niveau doctoral, Bac+6 et plus & Doctorat, PhD \\
\bottomrule
\end{tabularx}

Ainsi, un Master ou Bac+5 correspond au \textbf{niveau NCF 8}, une Licence au niveau 7,
un BTS/DUT au niveau 6 et un Doctorat au niveau 9. Pour les candidats, le diplôme déclaré,
plus précis, est utilisé en priorité ; le niveau d'étude général sert de solution de repli.

\begin{alertbox}
Le fichier de configuration classe actuellement \texttt{Master 1} au niveau 7, alors que la NCF
chargée dans le projet place Master I et II au niveau 8. Ce mapping doit être corrigé avant une
nouvelle exécution complète, sinon certaines compatibilités de qualification seront sous-estimées.
\end{alertbox}

\section{Le modèle a-t-il réellement été fine-tuné ?}

Oui. Un fine-tuning suppose que l'on parte d'un modèle préentraîné, que l'on calcule une perte
sur des données propres au domaine, puis que l'on modifie ses poids par rétropropagation.
Les artefacts du projet montrent ces trois éléments.

\begin{itemize}
  \item modèle de départ : \texttt{sentence-transformers/all-MiniLM-L6-v2} ;
  \item appel effectif à \texttt{trainer.train()} ;
  \item checkpoints contenant les poids, l'optimiseur et le scheduler ;
  \item cinq époques et 710 étapes enregistrées ;
  \item modèle final sauvegardé dans \texttt{models/st\_finetuned/final} ;
  \item perte observée passant d'environ 2,45 au début à environ 0,13 en fin d'entraînement.
\end{itemize}

\begin{oralbox}
« Il s'agit d'un vrai fine-tuning, car les poids du Transformer ont été optimisés sur des paires
du domaine emploi. La présence des checkpoints, des gradients, de l'état de l'optimiseur et du
modèle final exclut une simple réévaluation du modèle de base. »
\end{oralbox}

\newpage
\section{Protocole d'entraînement}

\begin{tabularx}{\textwidth}{>{\bfseries}p{0.34\textwidth}X}
\toprule
Élément & Valeur retenue \\
\midrule
Modèle initial & all-MiniLM-L6-v2 \\
Architecture & 6 couches Transformer, 12 têtes d'attention \\
Dimension des embeddings & 384 \\
Paires d'entraînement & 4\,542 \\
Paires de validation & 981 \\
Paires de test & 953 \\
Époques & 5 \\
Taille de lot & 32 \\
Taux d'apprentissage & $2\times10^{-5}$ \\
Perte & Multiple Negatives Ranking Loss \\
Scheduler & Cosine avec warmup \\
Weight decay & 0,01 \\
Critère du meilleur checkpoint & NDCG@10 de validation \\
Meilleur checkpoint & Époque 4, NDCG@10 = 0,6492 \\
Durée d'entraînement enregistrée & 15\,157 secondes, soit 4 h 12 min 37 s, sur CPU \\
\bottomrule
\end{tabularx}

La perte contrastive rapproche la représentation des métadonnées d'une offre de celle de sa
description. Dans un lot de 32 paires positives, les autres descriptions du lot servent de
négatifs implicites. Le modèle apprend donc à distinguer la bonne description des descriptions
concurrentes sans nécessiter une annotation négative manuelle pour chaque paire.

\begin{oralbox}
« Le fine-tuning complet sur cinq époques a duré 15\,157 secondes, soit environ 4 heures et
13 minutes sur CPU. Cette durée correspond à l'entraînement enregistré ; elle n'inclut pas toute
la préparation des données ni les analyses réalisées après l'entraînement. »
\end{oralbox}

\section{Quels paramètres ont été optimisés ?}

Le modèle final contient \textbf{22\,713\,216 paramètres}. L'audit du modèle montre que
\textbf{100\,\% de ces paramètres étaient entraînables}. Ils appartiennent à l'encodeur :

\begin{itemize}
  \item embeddings des tokens et des positions ;
  \item matrices Query, Key et Value de l'attention ;
  \item projections de sortie de l'attention ;
  \item couches feed-forward ;
  \item paramètres des LayerNorm des six couches Transformer.
\end{itemize}

Le mean pooling calcule la moyenne des représentations des tokens et la dernière étape normalise
le vecteur. Ces deux modules ne contiennent aucun poids entraînable. Le tokenizer et son
vocabulaire de 30\,522 tokens n'ont pas été réentraînés.

\begin{alertbox}
Les paramètres appris ne doivent pas être confondus avec les hyperparamètres. Les poids du
Transformer sont appris automatiquement ; le nombre d'époques, la taille de lot et le taux
d'apprentissage sont fixés avant l'entraînement.
\end{alertbox}

\newpage
\section{Pourquoi ce n'est pas du LoRA}

\begin{tabularx}{\textwidth}{>{\bfseries}p{0.25\textwidth}XX}
\toprule
Aspect & Fine-tuning réalisé & LoRA \\
\midrule
Poids initiaux & Tous entraînables & Généralement gelés \\
Paramètres ajoutés & Aucun adaptateur & Petites matrices de faible rang \\
Part entraînée & 100\,\% des 22,7 millions & Faible proportion des paramètres \\
Bibliothèque PEFT & Non utilisée & Habituellement utilisée \\
Sortie & Modèle complet sauvegardé & Adaptateurs séparés ou fusionnés \\
\bottomrule
\end{tabularx}

LoRA approxime la mise à jour d'une matrice $W$ par deux petites matrices :
\[
W' = W + BA,
\]
où le rang de $BA$ est faible. Cette stratégie réduit la mémoire nécessaire pour adapter les
grands modèles. Elle n'a pas été utilisée ici : aucune couche n'est gelée et aucun adaptateur
de faible rang n'est injecté.

\begin{oralbox}
« Nous avons effectué un fine-tuning complet, et non du LoRA. Ce choix reste réalisable parce
que MiniLM est un encodeur compact de 22,7 millions de paramètres, très inférieur à la taille
des grands modèles génératifs. »
\end{oralbox}

\section{Résultats avant et après fine-tuning}

\begin{tabularx}{\textwidth}{X>{\centering\arraybackslash}p{0.21\textwidth}>{\centering\arraybackslash}p{0.21\textwidth}>{\centering\arraybackslash}p{0.18\textwidth}}
\toprule
\textbf{Métrique} & \textbf{Modèle initial} & \textbf{Modèle fine-tuné} & \textbf{Gain absolu} \\
\midrule
NDCG@10 & 0,1681 & \textbf{0,6232} & +0,4552 \\
MRR@10 & 0,1410 & \textbf{0,5874} & +0,4463 \\
Recall@1 & 0,0986 & \textbf{0,5268} & +0,4281 \\
Recall@5 & 0,1983 & \textbf{0,6632} & +0,4648 \\
Recall@10 & 0,2560 & \textbf{0,7398} & +0,4837 \\
\bottomrule
\end{tabularx}

Le Recall@10 de 0,7398 signifie que le modèle retrouve environ 74\,\% des descriptions attendues
parmi les dix premiers résultats du protocole de test. Le NDCG@10 de 0,6232 indique que les bons
résultats sont aussi mieux placés en tête de liste qu'avec le modèle initial.

\begin{alertbox}
Cette évaluation mesure une tâche de récupération construite à partir des offres : retrouver la
description correspondant aux métadonnées de la même offre. Elle démontre une adaptation
sémantique au domaine, mais pas encore une pertinence candidat--offre validée par des recruteurs.
Les chiffres ne doivent donc pas être présentés comme une preuve directe de qualité du recrutement.
\end{alertbox}

\newpage
\section{Optimisation bayésienne avec l'échantillonneur TPE}

Le terme correct est \textbf{échantillonneur TPE}, et non \og échantillon TPE\fg{}. TPE signifie
\emph{Tree-structured Parzen Estimator}. Dans Optuna, cet algorithme choisit les paramètres à
tester au prochain essai en exploitant les résultats déjà observés.

\subsection{Principe général}

Après quelques essais initiaux, TPE sépare les combinaisons évaluées en deux groupes :

\begin{itemize}
  \item les essais ayant obtenu les meilleurs résultats ;
  \item les essais ayant produit des résultats moins bons.
\end{itemize}

Il estime ensuite deux distributions :
\[
l(w)=p(w\mid \text{bon résultat})
\qquad\text{et}\qquad
g(w)=p(w\mid \text{résultat moins bon}).
\]

TPE privilégie les paramètres pour lesquels le rapport suivant est élevé :
\[
\frac{l(w)}{g(w)}.
\]

Il recherche donc des valeurs fréquentes parmi les bons essais et rares parmi les mauvais.
L'algorithme ne calcule pas directement la solution optimale : il apprend progressivement dans
quelle zone effectuer les prochains essais, tout en conservant une part d'exploration.

\subsection{Application au score hybride}

Pour chaque couple candidat--offre, le score calibré prend la forme :
\[
S(c,o)=\alpha S_{\mathrm{sem}}(c,o)+\beta S_{\mathrm{Neo4j}}(c,o),
\qquad \alpha+\beta=1.
\]

À chaque essai, Optuna propose deux poids bruts, les normalise, recalcule le score des offres,
classe les résultats de chaque candidat et mesure le NDCG@10 moyen. Le processus réalisé dans
le projet comprend 500 essais sur 690 couples candidat--offre associés à 69 candidats, avec une
graine fixée à 42.

\begin{tabularx}{\textwidth}{>{\centering\arraybackslash}p{0.18\textwidth}>{\centering\arraybackslash}p{0.22\textwidth}X}
\toprule
\textbf{Poids Neo4j} & \textbf{NDCG@10} & \textbf{Lecture simplifiée} \\
\midrule
0,10 & 0,82 & Faible performance observée \\
0,30 & 0,87 & Amélioration du classement \\
0,70 & 0,94 & Zone prometteuse \\
0,80 & 0,96 & Bon essai ; TPE explorera probablement autour de cette zone \\
\bottomrule
\end{tabularx}

Cet exemple est pédagogique : ces quatre valeurs ne constituent pas les résultats bruts du projet.
Elles illustrent seulement la manière dont TPE déplace progressivement sa recherche.

\begin{alertbox}
L'optimisation utilise ici \texttt{proxy\_relevance} comme vérité de référence. Elle recherche donc
les poids qui reproduisent au mieux ce proxy, et non une pertinence certifiée par des recruteurs.
Le résultat doit être confirmé sur un jeu indépendant annoté par des experts. De plus, la solution
exacte \og Neo4j seul\fg{} a été évaluée comme cas limite après les essais TPE, car l'espace de
recherche imposait un poids strictement positif à chaque composante.
\end{alertbox}

\begin{oralbox}
« L'échantillonneur TPE est une méthode d'optimisation bayésienne qui modélise séparément les
bonnes et les mauvaises combinaisons de poids. Il propose ensuite les essais ayant la plus forte
probabilité d'améliorer le NDCG@10, sans devoir tester uniformément toutes les combinaisons. »
\end{oralbox}

\newpage
\section{Structure du graphe et dépendance aux hubs}

Cette analyse vérifie si le graphe peut relier candidats, compétences et offres sans dépendre
excessivement de quelques nœuds génériques. Elle intervient avant le reranking : un graphe trop
fragmenté explique mal les recommandations ; un graphe dominé par des hubs produit au contraire
des rapprochements trop larges.

\subsection{Comptages Cypher et projection matching}

Des requêtes Cypher comptent directement dans Neo4j les nœuds par label, les relations par type
et le nombre de voisins de chaque nœud. Le graphe complet contient 90\,230 nœuds et 550\,542
relations.

Une \textbf{projection matching} est ensuite construite en conservant les relations utiles :
\texttt{POSSEDE}, \texttt{REQUIERT}, \texttt{NECESSITE}, \texttt{CORRESPOND\_MEPC},
\texttt{A\_NIVEAU}, \texttt{REQUIERT\_NIVEAU\_NCF}, \texttt{DANS\_SECTEUR} et
\texttt{LOCALISEE\_A}. Cette projection compte 72\,516 nœuds et 499\,080 arêtes.

Pour étudier la connectivité, elle est transformée en graphe simple non dirigé avec NetworkX.
Le diagnostic vérifie alors l'existence d'une connexion sans interpréter son orientation.

\begin{alertbox}
La projection non dirigée convient au diagnostic de structure, mais pas à l'interprétation métier
finale : elle ignore le sens des relations et peut fusionner plusieurs liens entre les mêmes nœuds.
\end{alertbox}

\subsection{Degré et identification des hubs}

Le degré d'un nœud est son nombre de voisins. Les statistiques calculées par type de nœud sont la
moyenne, la médiane, le minimum, le maximum et la variance. Pour les compétences, le degré médian
est de 4, alors que le maximum atteint 6\,787. Les niveaux NCF atteignent jusqu'à 10\,261 connexions.

La distribution des degrés est représentée sur des axes logarithmiques afin de rendre visibles
simultanément les nombreux nœuds faiblement connectés et les quelques hubs. La longue traîne
observée signifie que la connectivité est très concentrée.

\begin{tabularx}{\textwidth}{>{\bfseries}p{0.27\textwidth}X}
\toprule
Indicateur & Interprétation \\
\midrule
Degré élevé & Le nœud est relié à de nombreuses entités ; il facilite la navigation. \\
Hub générique & Le nœud relie beaucoup d'objets mais discrimine peu les offres. \\
Degré médian faible & La majorité des nœuds possède seulement quelques relations ciblées. \\
Variance élevée & Les connexions sont très inégalement réparties entre les nœuds. \\
\bottomrule
\end{tabularx}

Un degré élevé n'est donc pas une preuve de pertinence. Une compétence générale comme la
communication peut être fréquente sans suffire à justifier une recommandation précise.

\subsection{Densité et composante géante}

Pour un graphe non dirigé comportant $n$ nœuds et $m$ arêtes, la densité est :
\[
D=\frac{2m}{n(n-1)}.
\]

Une faible densité n'est pas automatiquement un défaut : une offre doit être reliée uniquement
aux compétences, niveaux et secteurs qui la décrivent, et non à toutes les entités du graphe.

La composante géante est le plus grand groupe de nœuds accessibles les uns depuis les autres.
Dans la projection initiale, elle contient 72\,516 nœuds. Les principaux chemins
candidat--compétence--offre sont donc disponibles pour le GraphRAG.

\subsection{Test de robustesse par retrait ciblé}

Les nœuds sont classés par degré décroissant. Le programme retire ensuite les 1\,\%, 5\,\% et
10\,\% les plus connectés, puis recalcule la taille de la composante géante.

\begin{tabularx}{\textwidth}{>{\centering\arraybackslash}p{0.25\textwidth}>{\centering\arraybackslash}p{0.28\textwidth}X}
\toprule
\textbf{Hubs retirés} & \textbf{Composante géante restante} & \textbf{Lecture} \\
\midrule
0\,\% & 100\,\% & Projection initiale connectée \\
1\,\% & 30,1\,\% & Forte dépendance aux nœuds les plus connectés \\
5\,\% & 11,3\,\% & Fragmentation importante \\
10\,\% & 0,07\,\% & Effondrement presque total de la projection \\
\bottomrule
\end{tabularx}

Il s'agit d'une attaque ciblée et non d'une panne aléatoire. Le résultat montre une dépendance
structurelle aux hubs, mais ne prouve pas à lui seul que Neo4j est inutilisable. Le retrait touche
aussi des nœuds organisateurs comme les niveaux NCF, les secteurs et les localisations ; une forte
fragmentation est donc partiellement attendue.

\subsection{Protection contre le bruit de popularité}

Le score Neo4j n'utilise pas le degré brut. Il mesure la couverture des compétences :
\[
S_{\mathrm{Neo4j}}(c,o)=
\frac{|\mathcal{C}_{c}\cap\mathcal{C}_{o}|}{|\mathcal{C}_{o}|},
\]
où $\mathcal{C}_{c}$ désigne les compétences du candidat et $\mathcal{C}_{o}$ celles requises par
l'offre. Une compétence fréquente ne compte que si elle est possédée par le candidat et demandée
par l'offre. Cette règle limite le biais de popularité dans le reranking.

\begin{oralbox}
« L'analyse combine des comptages Cypher, une projection NetworkX, la distribution des degrés,
la composante géante et un retrait ciblé des hubs. Le graphe est connecté mais dépend de quelques
nœuds génériques. Nous utilisons donc la couverture des compétences, et non le degré brut, comme
signal Neo4j du score hybride. »
\end{oralbox}

\section{Trame recommandée pour le PowerPoint}

\begin{enumerate}
  \item \textbf{Diapositive 1 -- Pourquoi normaliser ?} Montrer un exemple brut puis normalisé.
  \item \textbf{Diapositive 2 -- Niveaux NCF.} Présenter uniquement les niveaux 5 à 9 les plus
  mobilisés dans les offres, avec Bac+5 $\rightarrow$ niveau 8.
  \item \textbf{Diapositive 3 -- Données de fine-tuning.} Afficher les trois tailles de split et
  la construction métadonnées $\rightarrow$ description.
  \item \textbf{Diapositive 4 -- Entraînement.} Présenter modèle initial, perte contrastive,
  cinq époques, batch 32 et meilleur checkpoint.
  \item \textbf{Diapositive 5 -- Paramètres entraînés.} Insister sur 22,7 millions de paramètres,
  100\,\% entraînables, et préciser « pas de LoRA ».
  \item \textbf{Diapositive 6 -- Résultats.} Comparer NDCG@10, MRR@10 et Recall@10 avant/après.
  \item \textbf{Diapositive 7 -- Calibration TPE.} Montrer la boucle proposition des poids,
  classement, calcul du NDCG@10 et nouvel essai.
  \item \textbf{Diapositive 8 -- Structure du graphe.} Présenter degré, hubs et composante géante.
  \item \textbf{Diapositive 9 -- Robustesse.} Montrer l'effet du retrait ciblé et expliquer
  pourquoi le score n'utilise pas le degré brut.
  \item \textbf{Diapositive 10 -- Limites.} Expliquer que le test valide le retrieval d'offres et
  que l'optimisation dépend d'un proxy ; une annotation humaine reste nécessaire.
\end{enumerate}

\section{Questions probables du jury}

\textbf{Pourquoi avoir fine-tuné MiniLM ?}\par
Le vocabulaire et les formulations des offres camerounaises diffèrent des textes généraux sur
lesquels le modèle initial a été préentraîné. L'adaptation contrastive améliore leur représentation.

\textbf{Pourquoi cinq époques ?}\par
Les métriques de validation progressent jusqu'à l'époque 4, puis reculent légèrement à l'époque 5.
Le mécanisme de meilleur checkpoint conserve donc l'état de l'époque 4. Les cinq époques servent
à observer cette stabilisation, pas à affirmer que la dernière époque est nécessairement la meilleure.

\textbf{Pourquoi ne pas utiliser LoRA ?}\par
Le modèle ne contient que 22,7 millions de paramètres. Un fine-tuning complet reste possible sur
la machine disponible et permet d'adapter directement tout l'encodeur.

\textbf{Le modèle apprend-il directement le matching candidat--offre ?}\par
Non. Il apprend d'abord une représentation adaptée aux offres. Le matching final combine ensuite
la similarité vectorielle et les relations du graphe. Une validation humaine indépendante reste
nécessaire pour mesurer la pertinence métier réelle.

\textbf{Comment TPE choisit-il le prochain essai ?}\par
Il compare les distributions des paramètres associés aux bons et aux moins bons résultats, puis
privilégie les valeurs fréquentes dans le premier groupe et rares dans le second. Il maintient
toutefois une part d'exploration afin de ne pas rester bloqué trop tôt dans une zone sous-optimale.

\textbf{Un hub très connecté est-il une compétence très pertinente ?}\par
Non. Le degré mesure la connectivité, pas la pertinence pour un candidat donné. Une compétence
générale peut être reliée à de nombreuses offres tout en étant peu discriminante. Le score utilise
donc la proportion des compétences requises effectivement possédées.

\section{Sources techniques vérifiées}

\begin{itemize}
  \item Configuration : \path{src/02_finetune_st/config_st.json}.
  \item Entraînement : \path{src/02_finetune_st/train_sentence_transformer.py}.
  \item Données : \path{data/finetune/pairs_train.jsonl}, \path{pairs_val.jsonl} et
  \path{pairs_test.jsonl}.
  \item État du trainer : \path{models/st_finetuned/checkpoints/checkpoint-710/trainer_state.json}.
  \item Modèle final : \path{models/st_finetuned/final/model.safetensors}.
  \item Comparaison : \path{models/st_finetuned/eval_comparatif/evaluation_comparatif.json}.
  \item Calibration TPE : \path{src/07_evaluation/optimize_hybrid_weights.py}.
  \item Résultats Optuna : \path{outputs/evaluation/hybrid_weight_optimization/best_weights_ndcg.json}.
  \item Diagnostic du graphe : \path{scripts/generate_neo4j_network_diagnostics.py}.
  \item Robustesse : dossier \path{outputs/memoire_stats/implementation_deep_current},
  fichier \path{18_neo4j_robustesse.csv}.
  \item Référentiel NCF : \path{data/processed/ncf_niveaux.parquet}.
  \item Référence scientifique : Reimers, N. et Gurevych, I. (2019),
  \emph{Sentence-BERT: Sentence Embeddings using Siamese BERT-Networks}, EMNLP-IJCNLP.
\end{itemize}

\end{document}
